% !TeX encoding = UTF-8
\documentclass{../plantilla/hvtbook}
\HVTTitulo{HidroEcoCaja}
\HVTSubtitulo{Sistema educativo y técnico para recuperar nuestra relación con el territorio y aprender a identificar, separar, registrar y aprovechar los residuos de una vivienda, empresa o finca. El jardín rural y la vermicultura unen saberes ancestrales y campesinos con observación técnica; la caracterización posterior permite definir el destino responsable de cada corriente.}
\HVTNivel{Proyecto de bioeconomía en Boyacá}
\HVTSerie{Publicación técnica}
\HVTNumero{04}
\HVTAccion{Explorar capítulos}
\HVTDescripcion{Publicación técnica del proyecto HidroEcoCaja para bioeconomía, aprovechamiento de residuos orgánicos y vermicultura en Boyacá.}
\HVTCanonical{https://hidrogenoverdeturquesa.com/proyecto-hidroecocaja}
\HVTRegreso{Volver a proyectos}{/\#portfolio}
\HVTPortada{../../images/portfolio/HidroEcoCaja.jpg}{Figura 1. Cuidar los residuos orgánicos también es cuidar el suelo, la memoria y el territorio.}
\begin{document}
\HVTMakeTitle

\begin{HVTPage}{Contenido}{Ruta del proyecto}{Del residuo separado a una decisión responsable}
\HVTLead{HidroEcoCaja no convierte automáticamente todos los residuos en hidrógeno. Enseña a reconocerlos, evitar su contaminación y producir información técnica para asignar a cada corriente el tratamiento que realmente le corresponde.}
\begin{HVTTaskOrdered}{Capítulos}
\item Propósito y alcance. \item Jardín rural, territorio y saberes. \item Diagnóstico de generación. \item Clasificación en el origen. \item Matriz de aceptación. \item Medición y trazabilidad. \item Preparación de la fracción orgánica. \item Diseño del módulo de vermicultura. \item Operación y seguimiento. \item Productos y desarrollo de biorremediadores. \item Vinculación con la cadena de valor del hidrógeno. \item Escalamiento en vivienda, empresa o finca. \item Seguridad, formación e indicadores.
\end{HVTTaskOrdered}
\end{HVTPage}

\begin{HVTPage}{Capítulo 1}{Propósito}{Aprender antes de transformar}
\HVTLead{El primer resultado del proyecto es una cultura de separación basada en evidencia. El usuario aprende qué genera, cuánto genera, qué puede recuperar y qué requiere un gestor especializado.}
\begin{HVTFlow}
\HVTFlowItem{1}{Reconocer}{Identificar el material, su origen y los posibles contaminantes.}
\HVTFlowItem{2}{Separar}{Evitar que las corrientes aprovechables pierdan calidad al mezclarse.}
\HVTFlowItem{3}{Decidir}{Asignar vermicultura, reciclaje, estudio técnico o gestión autorizada.}
\end{HVTFlow}
\begin{HVTWarning}{Principio del proyecto}Ningún residuo entra a un proceso por intuición. Su origen, composición, inocuidad y destino deben conocerse y registrarse.\end{HVTWarning}
\end{HVTPage}

\begin{HVTPage}{Capítulo 2}{Jardín rural}{Territorio, memoria y saberes}
\HVTLead{El jardín rural es un espacio vivo de aprendizaje donde se reconoce que la alimentación, el agua, el suelo, las plantas, los animales y los residuos forman parte de un mismo territorio. La HidroEcoCaja ayuda a recuperar esta relación al devolver atención y cuidado a la materia orgánica que normalmente se desecha.}
\begin{HVTFlow}
\HVTFlowItem{1}{Escuchar}{Reconocer conocimientos de familias, campesinos, mayores, custodios de semillas y comunidades sobre ciclos de siembra, suelos y plantas.}
\HVTFlowItem{2}{Practicar}{Clasificar residuos, cuidar las lombrices, observar la descomposición y relacionarla con el jardín y los alimentos.}
\HVTFlowItem{3}{Transmitir}{Registrar aprendizajes y compartirlos entre generaciones sin perder su procedencia ni reemplazar la experiencia local.}
\end{HVTFlow}
\end{HVTPage}

\begin{HVTPage}{Capítulo 2}{Jardín rural}{Aprendizaje territorial}
\begin{HVTTable}{Componente}{Aprendizaje territorial}{Actividad propuesta}
\HVTTableRow{Suelo}{Es un sistema vivo y no solamente un soporte}{Observar textura, cobertura, humedad, fauna y cambios en el tiempo}
\HVTTableRow{Semillas y plantas}{Expresan historia, alimentación y adaptación local}{Crear inventario participativo y calendario de siembra}
\HVTTableRow{Agua}{Conecta vivienda, cultivo, paisaje y salud}{Reconocer fuentes, recorridos, lluvias y prácticas de cuidado}
\HVTTableRow{Residuos orgánicos}{Pueden volver a ciclos biológicos si están bien separados}{Seguir su recorrido hasta la HidroEcoCaja}
\HVTTableRow{Memoria comunitaria}{Orienta prácticas adaptadas al lugar}{Conversatorios, recorridos y cuaderno de saberes}
\end{HVTTable}
\end{HVTPage}

\begin{HVTPage}{Capítulo 2}{Jardín rural}{Cuaderno y consentimiento}
\begin{HVTTask}{Cuaderno del jardín rural}
\item Nombre local de plantas, semillas, suelos y temporadas. \item Relatos sobre cultivo y aprovechamiento de residuos. \item Mapa del jardín, fuentes de agua y áreas de protección. \item Calendario de siembra, cosecha, poda y generación orgánica. \item Registro conjunto de observaciones tradicionales y mediciones técnicas.
\end{HVTTask}
\begin{HVTWarning}{Reconocimiento y consentimiento}Los saberes ancestrales, tradicionales y campesinos pertenecen a las personas y comunidades que los conservan. Su documentación o divulgación debe acordarse con ellas, reconocer su procedencia y respetar aquello que no deseen compartir.\end{HVTWarning}
\end{HVTPage}

\begin{HVTPage}{Capítulo 3}{Diagnóstico}{¿Qué residuos se generan?}
\HVTLead{La caracterización comienza donde se produce el residuo. Una vivienda, una oficina, un restaurante y una finca tienen perfiles distintos; por eso no se dimensiona una HidroEcoCaja con valores genéricos.}
\begin{HVTTable}{Origen}{Corrientes por revisar}{Condición que cambia la ruta}
\HVTTableRow{Vivienda}{Alimentos, jardín, empaques y residuos sanitarios}{Mezclas, aceites, medicamentos y limpieza}
\HVTTableRow{Empresa}{Comedor, cafetería, oficinas, producción y mantenimiento}{Actividad, sustancias usadas y requisitos del gestor}
\HVTTableRow{Finca}{Cosecha, poda, estiércoles, alimentos, envases e insumos}{Agroquímicos, medicamentos veterinarios y sanidad animal}
\end{HVTTable}
\HVTEquation{m_{\text{generada}}=\sum_i m_i}{}
\HVTText{Durante un periodo representativo se pesa cada corriente $i$, se registra su frecuencia y se anotan variaciones por temporada o actividad.}
\end{HVTPage}

\begin{HVTPage}{Capítulo 3}{Diagnóstico}{Entregables}
\begin{HVTTask}{Entregables}\item Mapa de puntos de generación. \item Inventario de recipientes y rutas internas. \item Registro de masa por tipo de residuo. \item Fotografías y observaciones de contaminación. \item Línea base antes de instalar el módulo.\end{HVTTask}
\end{HVTPage}

\begin{HVTPage}{Capítulo 4}{Clasificación}{Separación en el origen}
\HVTLead{En Colombia, la Resolución 2184 de 2019 unificó el código de colores: verde para orgánicos aprovechables, blanco para aprovechables y negro para no aprovechables. \HVTCite{1}}
\begin{HVTTable}{Color o corriente}{Contenido general}{Destino inicial}
\HVTTableRow{Verde}{Orgánicos aprovechables correctamente separados}{Evaluar vermicultura, compostaje u otra valorización}
\HVTTableRow{Blanco}{Plástico, vidrio, metales, papel, cartón y multicapa}{Entregar limpios y secos a reciclaje}
\HVTTableRow{Negro}{Residuos no aprovechables}{Servicio público o ruta local definida}
\HVTTableRow{Especiales}{Pilas, medicamentos, aparatos, aceites, químicos y peligrosos}{Posconsumo o gestor autorizado; nunca HidroEcoCaja}
\end{HVTTable}
\begin{HVTWarning}{La bolsa verde no es autorización automática}Antes de alimentar lombrices debe verificarse compatibilidad, ausencia de sustancias peligrosas y seguridad sanitaria.\end{HVTWarning}
\end{HVTPage}

\begin{HVTPage}{Capítulo 5}{Aceptación}{Qué entra y qué no entra}
\HVTLead{La lista definitiva dependerá de las pruebas del prototipo, la especie de lombriz, el entorno y el uso previsto del producto. Esta matriz conservadora orienta la fase inicial.}
\begin{HVTTable}{Categoría}{Ejemplos}{Acción}
\HVTTableRow{Aceptación inicial}{Restos vegetales, cáscaras, hojas secas y cartón limpio}{Trocear, mezclar y alimentar gradualmente}
\HVTTableRow{Evaluación previa}{Cítricos, cebolla, café, estiércoles, poda y grandes volúmenes}{Revisar origen, proporción, pretratamiento y respuesta}
\HVTTableRow{No admitir}{Plásticos, vidrio, metales, pañales, químicos, pilas, aceites y agroquímicos}{Separar y entregar por la ruta correspondiente}
\HVTTableRow{Riesgo sanitario}{Excretas humanas, animales enfermos, cadáveres, biosólidos o clínicos}{Gestión especializada; no usar en el módulo doméstico}
\end{HVTTable}
\begin{HVTWarning}{Diseño para bajos olores}La separación, aireación, humedad, carga gradual y protección permiten mantener bajos olores. Cambios de olor, moscas, líquido o calentamiento exigen revisión y corrección.\end{HVTWarning}
\end{HVTPage}

\begin{HVTPage}{Capítulo 6}{Trazabilidad}{Medir para aprender}
\HVTLead{Cada lote necesita identidad: fecha, procedencia, tipo de material, masa, observaciones y destino. Así se compara el desempeño y se evitan afirmaciones sin respaldo.}
\HVTEquation{T_{\text{separación}}=\frac{m_{\text{correctamente separada}}}{m_{\text{generada}}}\times 100}{}
\begin{HVTTable}{Dato}{Instrumento o registro}{Uso}
\HVTTableRow{Masa de entrada y rechazo}{Balanza y formato por lote}{Capacidad y eficiencia de separación}
\HVTTableRow{Temperatura}{Termómetro de sonda}{Detectar calentamiento anormal}
\HVTTableRow{Humedad}{Observación estandarizada o sensor validado}{Evitar sequedad o saturación}
\HVTTableRow{pH, si se requiere}{Medidor calibrado}{Comparar mezclas y diagnosticar cambios}
\HVTTableRow{Olor, insectos y comportamiento}{Lista de verificación}{Acción correctiva temprana}
\end{HVTTable}
\end{HVTPage}

\begin{HVTPage}{Capítulo 6}{Trazabilidad}{Ficha mínima de lote}
\begin{HVTTask}{Ficha mínima de lote}\item Código y responsable. \item Origen y descripción verificable. \item Masa incorporada y materiales de estructura. \item Fecha de alimentación y controles. \item Incidentes, correcciones y destino final.\end{HVTTask}
\end{HVTPage}

\begin{HVTPage}{Capítulo 7}{Preparación}{Acondicionar la fracción orgánica}
\HVTLead{Las lombrices y los microorganismos actúan sobre una mezcla preparada, no sobre una bolsa heterogénea. El acondicionamiento busca alimentación homogénea, aireada y libre de impropios.}
\begin{HVTFlow}\HVTFlowItem{1}{Inspeccionar}{Retirar empaques, etiquetas, grapas, vidrio y materiales extraños.}\HVTFlowItem{2}{Preparar}{Reducir tamaño, mezclar materiales húmedos y estructurantes y preacondicionar.}\HVTFlowItem{3}{Probar}{Incorporar una cantidad pequeña y observar antes de aumentar la carga.}\end{HVTFlow}
\begin{HVTWarning}{Parámetros por validar}No se fijará una receta universal de humedad, pH, relación carbono/nitrógeno ni residencia. Los rangos se establecerán con bibliografía, condiciones locales y ensayos documentados. \HVTCite{3}\end{HVTWarning}
\end{HVTPage}

\begin{HVTPage}{Capítulo 8}{Diseño}{Módulo de vermicultura}
\HVTLead{El diseño debe proteger a los organismos, facilitar el aprendizaje y permitir inspección, limpieza y mantenimiento sin exponer al usuario a partes inseguras.}
\begin{HVTTable}{Subsistema}{Función}{Criterio de diseño}
\HVTTableRow{Estructura}{Soportar carga y separar niveles}{Estabilidad, durabilidad y compatibilidad}
\HVTTableRow{Lecho}{Alojar lombrices y mezcla}{Volumen útil, acceso y protección de luz}
\HVTTableRow{Ventilación}{Mantener condiciones aerobias}{Aberturas protegidas contra plagas}
\HVTTableRow{Control de agua}{Evitar inundación y derrames}{Cubierta, drenaje y contención secundaria}
\HVTTableRow{Cosecha}{Separar material procesado}{Extracción sin mezclar lotes ni dañar el sistema}
\HVTTableRow{Instrumentación}{Observar variables}{Puntos accesibles y repetibles}
\end{HVTTable}
\HVTEquation{M_{\text{en proceso}}\approx\dot{m}_{\text{aceptada}}\,t_{\text{residencia}}}{}
\HVTText{Es una primera aproximación. El volumen real incluye densidad aparente, lecho, aireación, variaciones de carga y margen operativo.}
\end{HVTPage}

\begin{HVTPage}{Capítulo 9}{Operación}{Rutina de seguimiento}
\HVTLead{La alimentación se incrementa únicamente cuando el sistema responde de forma estable. El registro relaciona causas probables con acciones correctivas.}
\begin{HVTTable}{Señal}{Revisión inmediata}{Respuesta inicial}
\HVTTableRow{Olor intenso}{Exceso, saturación o compactación}{Suspender carga y recuperar aireación}
\HVTTableRow{Acumulación de líquido}{Lluvia, humedad o drenaje bloqueado}{Contener, corregir y no aplicar sin evaluación}
\HVTTableRow{Calentamiento}{Carga fresca o pretratamiento insuficiente}{Detener alimentación y proteger lombrices}
\HVTTableRow{Fuga o mortalidad}{Temperatura, humedad, irritantes o falta de alimento}{Aislar lote y revisar historial}
\HVTTableRow{Plagas}{Alimentos expuestos, tapas y limpieza}{Cubrir y retirar la causa}
\end{HVTTable}
\begin{HVTTask}{Plan de operación}\item Responsable y frecuencia de inspección. \item Procedimiento de alimentación y cosecha. \item Limpieza compatible. \item Respuesta a derrames, plagas y fallas. \item Protección frente a lluvia, sol, animales y manipulación.\end{HVTTask}
\end{HVTPage}

\begin{HVTPage}{Capítulo 10}{Biorremediadores}{De la HidroEcoCaja al biorremediador}
\HVTLead{La HidroEcoCaja genera productos biológicos de base mediante vermicultura controlada. Tras su cosecha pasan a un procedimiento independiente de selección, acondicionamiento o formulación para una aplicación concreta.}
\HVTText{La literatura documenta el uso del vermicompost y enmiendas microbianas en biorremediación. El procedimiento debe conservar trazabilidad, definir el contaminante objetivo y comprobar eficacia. \HVTCite{6}}
\begin{HVTTable}{Salida}{Verificación}{Restricción}
\HVTTableRow{Vermicompost}{Homogeneidad, estabilidad y análisis}{No atribuir propiedades sin respaldo}
\HVTTableRow{Material no procesado}{Separación y causa del rechazo}{Reprocesar solo si el protocolo lo permite}
\HVTTableRow{Líquido de drenaje}{Volumen, origen y caracterización}{No llamarlo fertilizante ni aplicarlo automáticamente}
\HVTTableRow{Datos del proceso}{Integridad de registros y balance}{No sustituye ensayos de laboratorio}
\end{HVTTable}
\end{HVTPage}

\begin{HVTPage}{Capítulo 10}{Biorremediadores}{Secuencia de desarrollo}
\begin{HVTFlow}\HVTFlowItem{1}{Obtener la base}{Cosechar productos trazables de materias primas conocidas y operación controlada.}\HVTFlowItem{2}{Formular}{Aplicar selección, acondicionamiento, combinación o activación definida.}\HVTFlowItem{3}{Comprobar}{Caracterizar y validar la respuesta frente al contaminante y suelo específicos.}\end{HVTFlow}
\begin{HVTTable}{Línea}{Pregunta experimental}{Evidencia requerida}
\HVTTableRow{Enmienda orgánica}{¿Mejora condiciones del suelo?}{Caracterización, dosis, control y seguimiento}
\HVTTableRow{Bioestimulación}{¿Estimula microorganismos nativos?}{Cambio del contaminante y actividad biológica}
\HVTTableRow{Bioaumentación}{¿Aporta una función reproducible?}{Identificación, concentración, pureza y eficacia}
\HVTTableRow{Inmovilización}{¿Disminuye movilidad o biodisponibilidad?}{Ensayos químicos y ecotoxicológicos}
\end{HVTTable}
\end{HVTPage}

\begin{HVTPage}{Capítulo 10}{Biorremediadores}{Validación y requisitos}
\begin{HVTTaskOrdered}{Protocolo mínimo de validación}\item Definir contaminante y objetivo. \item Caracterizar suelo y producto. \item Establecer tratamientos, control y réplicas. \item Medir concentración, toxicidad o biodisponibilidad. \item Evaluar inocuidad, estabilidad y efectos secundarios. \item Documentar alcance y condiciones.\end{HVTTaskOrdered}
\begin{HVTWarning}{Secuencia de desarrollo}HidroEcoCaja produce materiales de base; un procedimiento posterior los transforma en biorremediadores. Para usarlos o comercializarlos como bioinsumos debe demostrarse su función y cumplirse la Resolución ICA 68370 de 2020. \HVTCite{7}\end{HVTWarning}
\end{HVTPage}

\begin{HVTPage}{Capítulo 11}{Cadena H2}{De la clasificación a una ruta tecnológica}
\HVTLead{Biomasa y aguas residuales pueden estudiarse en procesos microbianos; otros orgánicos, en rutas termoquímicas. Exigen caracterización, pretratamiento, reactor, control, seguridad y validación. Son procesos distintos de la vermicultura. \HVTCite{4} \HVTCite{5}}
\begin{HVTFlow}\HVTFlowItem{1}{Corriente limpia}{Separación, pesaje, origen conocido y ausencia de contaminantes.}\HVTFlowItem{2}{Caracterización}{Composición, humedad, biodegradabilidad, cenizas y variabilidad.}\HVTFlowItem{3}{Ruta validada}{Vermicultura, reciclaje, biohidrógeno, ruta termoquímica o gestor.}\end{HVTFlow}
\begin{HVTTable}{Corriente}{Ruta que puede estudiarse}{Decisión necesaria}
\HVTTableRow{Orgánicos húmedos}{Vermicultura, compostaje o investigación microbiana}{Inocuidad, biodegradabilidad y desempeño}
\HVTTableRow{Biomasa seca}{Aprovechamiento material o termoquímica}{Humedad, composición, energía y emisiones}
\HVTTableRow{Reciclables limpios}{Reciclaje}{No desviarlos sin justificación de ciclo de vida}
\HVTTableRow{Peligrosos}{Gestión especializada}{Excluir del prototipo}
\end{HVTTable}
\begin{HVTWarning}{Puerta de decisión}La HidroEcoCaja prepara capacidades y datos; no es un reactor de hidrógeno ni garantiza viabilidad técnica, ambiental o económica.\end{HVTWarning}
\end{HVTPage}

\begin{HVTPage}{Capítulo 12}{Escalamiento}{Una metodología, tres contextos}
\HVTLead{La secuencia se conserva, pero cambian capacidad, responsables, controles y obligaciones.}
\begin{HVTTable}{Aplicación}{Enfoque}{Requisitos adicionales}
\HVTTableRow{Vivienda}{Formación familiar y pequeña generación}{Operación simple, higiene, niños y mascotas}
\HVTTableRow{Empresa}{Separación por áreas y trazabilidad}{Responsables, capacitación, almacenamiento y gestores}
\HVTTableRow{Finca}{Residuos domésticos, agrícolas y pecuarios}{Bioseguridad, agroquímicos y sanidad animal}
\end{HVTTable}
\begin{HVTTaskOrdered}{Fases del prototipo}\item Diagnóstico y diseño conceptual. \item Construcción instrumentada. \item Puesta en marcha con bajo riesgo. \item Validación de capacidad y protocolo. \item Replicación modular documentada.\end{HVTTaskOrdered}
\HVTEquation{N_{\text{módulos}}=\left\lceil\frac{\dot{m}_{\text{aceptada total}}}{\dot{m}_{\text{validada por módulo}}}\right\rceil}{}
\HVTText{La capacidad por módulo debe provenir de ensayos propios; no de cifras comerciales sin condiciones de prueba.}
\end{HVTPage}

\begin{HVTPage}{Capítulo 13}{Gestión}{Seguridad, formación e indicadores}
\HVTLead{El proyecto se completa cuando las personas pueden operar el módulo, explicar el destino de cada residuo y demostrar resultados con registros.}
\begin{HVTTable}{Indicador}{Cálculo o unidad}{Qué demuestra}
\HVTTableRow{Generación total}{kg/periodo}{Línea base y tendencias}
\HVTTableRow{Separación correcta}{\% de masa generada}{Calidad del aprendizaje}
\HVTTableRow{Orgánico aceptado}{kg/periodo}{Carga real del sistema}
\HVTTableRow{Rechazos e impropios}{kg y tipo}{Errores por corregir}
\HVTTableRow{Incidentes}{número, causa y cierre}{Seguridad y mejora continua}
\HVTTableRow{Producto caracterizado}{kg por lote}{Resultado verificable}
\end{HVTTable}
\end{HVTPage}

\begin{HVTPage}{Capítulo 13}{Gestión}{Entregable integral}
\begin{HVTTask}{Entregable integral}\item Manual ilustrado de separación y aceptación. \item Planos, materiales y memoria de diseño. \item Protocolo de operación y contingencias. \item Formatos y tablero de indicadores. \item Plan de formación. \item Informe de validación y escalamiento.\end{HVTTask}
\end{HVTPage}

\begin{HVTPage}{Fuentes}{Base técnica}{Referencias consultadas}
\begin{HVTReferences}
\HVTReference{Ministerio de Ambiente. \emph{Código de colores para separación en la fuente, Resolución 2184 de 2019}}{https://www.minambiente.gov.co/gobierno-unifica-el-codigo-de-colores-para-la-separacion-de-residuos-en-la-fuente-a-nivel-nacional/}
\HVTReference{United States Environmental Protection Agency. \emph{Composting}}{https://www.epa.gov/sustainable-management-food/composting}
\HVTReference{Thirunavukkarasu, A. et al. (2023). \emph{Sustainable organic waste management using vermicomposting}}{https://doi.org/10.1039/D2EM00324D}
\HVTReference{U.S. Department of Energy. \emph{Hydrogen Production: Microbial Biomass Conversion}}{https://www.energy.gov/cmei/fuels/hydrogen-production-microbial-biomass-conversion}
\end{HVTReferences}
\end{HVTPage}
\begin{HVTPage}{Fuentes}{Base técnica}{Referencias consultadas}
\begin{HVTReferences}
\HVTReference{U.S. Department of Energy. \emph{Hydrogen Production: Biomass Gasification}}{https://www.energy.gov/cmei/fuels/hydrogen-production-biomass-gasification}
\HVTReference{Kaur, T. (2022). \emph{Vermicomposting with microbial amendment: implications for bioremediation}}{https://doi.org/10.5114/bta.2022.116213}
\HVTReference{Instituto Colombiano Agropecuario. \emph{Resolución 68370 de 2020}}{https://www.ica.gov.co/normatividad/normas-ica/resoluciones-oficinas-nacionales/2020/2020r68370}
\end{HVTReferences}
\HVTText{Las referencias orientan el diseño conceptual. Los rangos de operación, la inocuidad, el desempeño y la viabilidad deben validarse para residuos, clima y escala reales. Esta publicación contiene únicamente información autorizada; los resultados detallados y estratégicos son confidenciales.}
\end{HVTPage}

\begin{HVTPage}{Colaboración}{Conversemos}{¿Necesitas investigar, formular o evaluar un proyecto relacionado?}
\HVTLead{Conversemos sobre tu necesidad de bioeconomía, aprovechamiento de residuos orgánicos, vermicultura o diseño de una solución relacionada.}
\HVTContact{Consultar proyecto de bioeconomía por WhatsApp}{https://wa.me/573209574884}
\end{HVTPage}
\end{document}
